\documentclass[11pt]{amsart}
\usepackage[T1]{fontenc}
\usepackage{lmodern}
\usepackage{amsmath,amssymb,amsthm,mathtools}
\usepackage{booktabs}
\usepackage{microtype}
\usepackage{xcolor}
\usepackage[colorlinks=true,linkcolor=blue!55!black,citecolor=blue!55!black,urlcolor=blue!55!black]{hyperref}

\newtheorem{theorem}{Theorem}[section]
\newtheorem{lemma}[theorem]{Lemma}
\newtheorem{proposition}[theorem]{Proposition}
\theoremstyle{remark}
\newtheorem{remark}[theorem]{Remark}

\newcommand{\PP}{\mathbb P}

\title{The Prime-Cube/Prime-Square Minimum Below One Trillion Is 49,600}
\author{Zijian Zeng}
\address{Institute of Computer Science and Digital Innovation, UCSI University, Kuala Lumpur, 56000, MALAYSIA}
\email{zijianzeng@foxmail.com}
\email{1002266693@ucsiuniversity.edu.my}
\date{August 31, 2026}
\subjclass[2020]{Primary 11Y50; Secondary 11G05, 11A41}
\keywords{Mordell equation, integral point, prime square, prime cube, exact computation}
\hypersetup{%
  pdftitle={The Prime-Cube/Prime-Square Minimum Below One Trillion Is 49,600},
  pdfauthor={Zijian Zeng}
}

\begin{document}

\begin{abstract}
For every prime $p$ with $10^6\le p\le10^{12}$, consider the squares of the
two primes adjacent to $p^{3/2}$.  We prove that the least distance from
$p^3$ to either such square is exactly $49{,}600$, attained at
$p=1{,}587{,}809$ and $q=2{,}000{,}771{,}023$.  The main observation is that
every possible improvement is an integral point on a Mordell curve
$Y^2=X^3+k$ with $0<|k|<49{,}600$.  A published complete classification for
$0<|k|\le10^7$, followed by an independently checked exact prime filter,
excludes all improvements.  Exact primality certificates and an integer
square-root computation certify attainment.
\end{abstract}

\maketitle

\section{Introduction}

For a prime $p$, let $q_-(p)<p^{3/2}<q_+(p)$ denote the primes adjacent to
$p^{3/2}$.  The problem recorded as TheoremDB P2484 asks for the exact value
of
\begin{equation}\label{eq:delta}
 \Delta=\min_{\substack{p\in\PP\\10^6\le p\le10^{12}}}
 \min\{p^3-q_-(p)^2,\ q_+(p)^2-p^3\}.
\end{equation}
The July 24, 2026 snapshot reports an exact sweep through $p=10^{11}$, with
incumbent $49{,}600$, while leaving the final interval open
\cite{TheoremDBP2484}.  The incumbent primality and adjacency calculation is
also recorded separately \cite{TheoremDBR560}.

The unexploited feature is that a putative improvement produces an integral
point on one of finitely many Mordell curves.  Bennett and Ghadermarzi gave a
complete resolution of
\begin{equation}\label{eq:mordell}
 Y^2=X^3+k
\end{equation}
for every nonzero integer $k$ with $|k|\le10^7$ and posted their full output
\cite{BennettGhadermarzi2015}.  Their range is more than sufficient for
\eqref{eq:delta}.  Filtering those published tables yields the following
exact resolution.

\begin{theorem}\label{thm:main}
The value in \eqref{eq:delta} is
\[
 \Delta=49{,}600.
\]
Equality is attained for
\[
 p=1{,}587{,}809,\qquad q_-(p)=2{,}000{,}771{,}023.
\]
\end{theorem}

Searches through August 31, 2026, using the exact attaining pair, the three
smaller integer-square candidates, and the problem title located the
TheoremDB prefix computation and the Mordell tables, but no source spelling
out Theorem~\ref{thm:main}.  This is only a bounded literature-search
statement; the mathematical result below does not depend on a priority
claim.

\section{Reduction to Mordell equations}

\begin{lemma}\label{lem:adjacent}
For every prime $p$ in the range of \eqref{eq:delta},
\[
 \min\{p^3-q_-(p)^2,q_+(p)^2-p^3\}
 =\min_{q\in\PP}|p^3-q^2|.
\]
Consequently,
\begin{equation}\label{eq:twoprime}
 \Delta=\min\{|p^3-q^2|:p,q\in\PP,\ 10^6\le p\le10^{12}\}.
\end{equation}
\end{lemma}

\begin{proof}
The map $q\mapsto q^2$ is strictly increasing for positive integers.  Among
primes below $p^{3/2}$, the largest one therefore gives the closest square
from below; among primes above, the smallest gives the closest square from
above.  Taking the smaller of these two distances proves the assertion.
\end{proof}

The difference in \eqref{eq:twoprime} is never zero.  Indeed, if
$p^3=q^2$ for primes $p,q$, unique factorization would require the same prime
to occur to both exponent three and exponent two.

For a nonzero difference, set
\[
 k=q^2-p^3.
\]
Then $(X,Y)=(p,q)$ is an integral solution of \eqref{eq:mordell}.  Thus every
pair that could improve the incumbent has $0<|k|<49{,}600$ and is covered by
the Bennett--Ghadermarzi classification.

\section{The exact filter}

We used the two complete tables for negative and positive $k$ with
$0<|k|\le10^6$ linked from the authors' data page
\cite{BennettGhadermarzi2015}.  This smaller portion of their published
range already contains every curve needed here.

\begin{proposition}\label{prop:filter}
In the published tables, there are $65{,}522$ signed integral points with
$0<|k|\le49{,}600$.  Exactly $196$ of them have
$10^6\le X\le10^{12}$, or $98$ after identifying $Y$ and $-Y$.
Among these, only the following four have prime $X$.
\begin{center}
\begin{tabular}{rrrl}
\toprule
$k$ & $X$ & $|Y|$ & classification \\
\midrule
$ 17{,}767$ & $15{,}175{,}973$ & $59{,}120{,}053{,}422$ & $Y$ even \\
$-24{,}315$ & $ 1{,}646{,}719$ & $ 2{,}113{,}144{,}738$ & $Y$ even \\
$ 25{,}895$ & $103{,}289{,}609$ & $1{,}049{,}747{,}744{,}368$ & $Y$ even \\
$-49{,}600$ & $ 1{,}587{,}809$ & $ 2{,}000{,}771{,}023$ & both prime \\
\bottomrule
\end{tabular}
\end{center}
\end{proposition}

\begin{proof}
The accompanying discovery program parses every table entry, checks
$Y^2=X^3+k$ in integer arithmetic, applies the two bounds, identifies the
two signs of $Y$, and classifies the remaining coordinates.  A separate
standard-library checker starts again from the two source PDFs, uses a
different text-extraction mode and parser, recomputes every displayed
Mordell identity, and uses exact trial division for the $98$ relevant
points.  Both computations give the displayed counts and rows.
\end{proof}

The proposition is an exact filter of a cited complete classification.  It
does not, by itself, reprove that the authors' tables contain every integral
point; that completeness is the imported result of
\cite{BennettGhadermarzi2015}.

\section{Primality, adjacency, and proof of the main theorem}

For completeness, the attaining pair has short $n-1$ certificates.  The
complete factorizations are
\begin{align*}
 1{,}587{,}809-1&=2^5\cdot29^2\cdot59,\\
 2{,}000{,}771{,}023-1&=2\cdot3\cdot19^2\cdot337\cdot2741.
\end{align*}
All factors on the right are prime by trial division.  For each distinct
prime factor $s$ of $n-1$, the following table gives a base $a$ for which
\[
 a^{n-1}\equiv1\pmod n,
 \qquad
 \gcd\bigl(a^{(n-1)/s}-1,n\bigr)=1.
\]
\begin{center}
\begin{tabular}{rll}
\toprule
$n$ & prime factors $s$ & corresponding bases $a$ \\
\midrule
$1{,}587{,}809$ & $2,29,59$ & $3,2,2$ \\
$2{,}000{,}771{,}023$ & $2,3,19,337,2741$ & $3,11,2,2,2$ \\
\bottomrule
\end{tabular}
\end{center}
Pocklington's criterion \cite{Pocklington1914} therefore proves that both
integers are prime.  These data agree with the earlier incumbent certificate
\cite{TheoremDBR560} and are rechecked by the independent verifier.

Exact arithmetic gives
\begin{align*}
 1{,}587{,}809^3&=4{,}003{,}084{,}686{,}476{,}516{,}129,\\
 2{,}000{,}771{,}023^2&=4{,}003{,}084{,}686{,}476{,}466{,}529,
\end{align*}
so their difference is $49{,}600$.  Moreover,
\[
 \left\lfloor\sqrt{1{,}587{,}809^3}\right\rfloor
 =2{,}000{,}771{,}023.
\]
Since this integer is prime, it is the adjacent prime below
$1{,}587{,}809^{3/2}$.

\begin{proof}[Proof of Theorem~\ref{thm:main}]
The certified attaining pair gives $\Delta\le49{,}600$.  If a smaller pair
existed, Lemma~\ref{lem:adjacent} would give primes $p,q$ in the required
$p$-range with $0<|p^3-q^2|<49{,}600$.  Setting $k=q^2-p^3$ would produce an
integral point covered by the complete Bennett--Ghadermarzi tables.
Proposition~\ref{prop:filter} shows that no such point has both $X$ and $|Y|$
prime.  Hence $\Delta\ge49{,}600$, proving equality.
\end{proof}

\section{Reproducibility and limitations}

The ancillary files contain the source-table locators, a discovery program,
the exact filtered certificate, the Pocklington data, and an independent
checker.  The checker uses CPython's standard library and Poppler's text
extractor.  A single documented command reruns all identity, range,
primality, and adjacency checks.  Integrity information and the complete
file inventory are recorded in a separate release manifest rather than in
this article.

The only nonlocal completeness dependency is the published
Bennett--Ghadermarzi computation.  The ancillary checker verifies the posted
data and the deduction made here, but does not rerun their month-long
enumeration of cubic forms and Thue equations.  No proof assistant was used.

\section*{Data availability}

The integral-point tables are available from the data page accompanying
\cite{BennettGhadermarzi2015}.  The derived certificate and verification
software are included with this manuscript.

\bibliographystyle{amsplain}
\bibliography{references}

\end{document}
